\section{The half-plane coboundary}
\label{sec:coboundary}

Let $m=\ind_M$ for the tie-symmetrized selector of
Section~\ref{sec:seed}.  We use
\[
  \ind_{\RC(M)}(w):=m(\RC(w)),
\]
which agrees with set-image notation because reverse complement is an
involution.

\begin{theorem}[Exact reverse-complement defect]
\label{thm:coboundary}
For every length-$k$ word $w$,
\begin{equation}
  m(\RC w)-m(w)
   =\chi(t(w))-\chi(h(w))
   =(\nabla\chi)(w).
  \label{eq:coboundary}
\end{equation}
Equivalently, $\ind_{\RC(M)}-\ind_M=\nabla\chi$.
\end{theorem}

\begin{proof}
Set
\[
  A=\operatorname{Im}P(w),
  \qquad B=\operatorname{Im}P(Rw).
\]
Lemma~\ref{lem:covariance} sends the ordered imaginary coordinates of
$\RC(w)$ to $(B,A)$.  If $A$ and $B$ are nonzero, the selector rule gives
\[
  m(\RC w)-m(w)
   =\ind[B<0<A]-\ind[A<0<B]
   =\ind[A>0]-\ind[B>0].
\]

It remains to check the rays.  If $A=0$, write $P(w)=r\in\R$.  Then
$B=-r\sin(2\pi/k)$.  For $r<0$ both sides are $-1$, while for $r>0$ both
are zero.  If $B=0$, write $P(Rw)=r\in\R$.  For $r>0$ both sides are $+1$,
while for $r<0$ both are zero.  When $k>2$, $A=B=0$ with $P(w)\neq0$ is
impossible because $P(Rw)=\zeta^{-1}P(w)$ and $\zeta$ is not real.  In the
zero sector, the chosen ties are RC invariant, so the left side vanishes;
both lower imaginary coordinates vanish as well.

Finally Lemma~\ref{lem:covariance} identifies the last sign difference with
$\chi(t(w))-\chi(h(w))$, with the gradient convention
\eqref{eq:gradient}.
\end{proof}

Equation~\eqref{eq:coboundary} is the bridge between the cyclotomic seed and
an exact legal move path.  Its right-hand side is not a heuristic score: it
is the literal indicator defect, edge by edge.  The use of an \emph{open}
upper half-plane and the chosen negative-ray convention are load-bearing.
